\chapter{The Whole Machine}
% full version: chapters/10_synthesis.tex

\pencilfig{10}{\pencilart{10}}{The whole machine: the data bus, the brakes in it, and four radio control channels above.}

We took one type of neuron at a time and asked what it adds to the machine. It is time to put the parts together.

The result is a machine of three layers. At the bottom is the huge telephone network of \nt{GLUT} neurons: the data run through it. Next to it are the braking \nt{GABA} neurons: they control the flow and decide what gets through, when, and how loudly. And above it all are a few radio stations, each with its own knob. \nt{DOPA} says which changes in the connections to keep. \nt{SERO} holds the overall level and, according to the hypothesis, the horizon of patience. \nt{NORA} chooses between exploring and deciding. \nt{ACET}, between writing and reading. Each radio station is really a small bundle of channels, not one wire, but there is still only a handful of channels for eighty-six billion neurons.

\section*{One Story}

Let us follow one event through the whole machine. An animal has long known: press lever A and you get juice. One day the world changes: A no longer works, and the juice now comes from lever B, which the animal has hardly used.

No juice comes after A: dopamine answers with a dip, and the connections that chose A weaken. The dips repeat, the errors become larger than usual, and noradrenaline, according to the hypothesis, reports: the world has changed. The contrast knob goes down, the choice becomes softer, and the noise more often nudges the animal to try B. Acetylcholine switches the network into write mode. Trying B brings unexpected juice: dopamine flares, and the connections that chose B grow stronger. After a few attempts the expectation catches up with the new world, the surprises end, and the knobs return to their places.

Note that no knob knows what the others are doing. Each responds to its own cues, and the sequence of actions assembles itself, as in an asynchronous circuit where each block fires when its inputs are ready. That each knob behaves this way on its own is mostly measured. That they work together in concert is still a hypothesis.

\section*{How It Differs from Your Laptop}

A laptop has a common clock; the brain does not. The elements of a laptop are reliable; a synapse loses more than half of its signals. A laptop's memory is separate from its processor; in the brain it sits in the same connections. A laptop learns (if it learns) by backpropagation of error; the brain, by local rules and a single broadcast ``better than expected.'' And most important: a laptop spends tens of watts on one processor, while the brain spends twenty watts on everything.

\section*{What We Don't Know}

We do not know how much the cortex uses the precise timing of clicks. We do not know how it tunes billions of connections in many-storied circuits when it has one teacher for all. We do not fully understand what the radio stations transmit. We do not know how distant parts of the brain coordinate their work with no common clock. And so far nobody knows how to build a machine that would do the same on twenty watts.

From here the book goes beyond this core: to the inputs and outputs of the machine, to debugging it, to sleep, and to machines made of living neurons on a chip.

\fullversion{10}
